\documentclass{scrartcl}
\usepackage[utf8]{inputenc}
\usepackage{hyperref}
\usepackage{url}
\usepackage{natbib}
\usepackage{graphicx}
\usepackage{cleveref} % this must be the last package to be loaded

\newcommand{\emailaddr}[1]{\href{mailto:#1}{\texttt{#1}}}

\title{\LARGE
    Final Report:\\
    Orbital
}
\subtitle{Cryptocurrency Market Data \& News Aggregation Platform: Distributed Systems Course\\
2025--2026}

\author{
    Aldo Visconti \\ \emailaddr{aldo.visconti@studio.unibo.it}
}

\date{\today}

\begin{document}

\maketitle

\noindent
\textbf{Orbital} is a distributed, microservices-based platform that aggregates real-time
cryptocurrency market data and financial news behind a single API gateway, with a Kotlin
Multiplatform client (desktop and web) on top. Application logic is decomposed into two
specialized services, \textbf{market-service} and \textbf{news-service}, fronted by a
\textbf{gateway} that proxies, caches, and rate-limits every client request.

The project focuses on keeping that aggregated data available and durable independently of any
single component failing. The \textbf{CAP theorem} trade-off is explicitly addressed: the system
prioritizes \textbf{availability} over strict consistency, since a client is always better served
a slightly stale cached price than an error. \textbf{Fault tolerance} is implemented through
real, verified infrastructure: both PostgreSQL databases run with a streaming-replication
standby, and the Redis cache runs under a three-node Sentinel quorum, so a single-instance
failure does not take the gateway's caching or rate-limiting down with it. Market data is
continuously ingested by a standalone worker, decoupled from API traffic, which persists an
append-only time series and performs a one-time historical backfill for newly observed coins.

\section*{Disclaimer}

During the preparation of this work, the author used Claude Code (Anthropic) as a pair-programming
and writing assistant: for implementing and debugging backend services, designing and drafting the
Postgres replication / Redis Sentinel / distributed rate-limiting infrastructure, performing a
security and code-quality audit of the resulting system, and drafting portions of this report from
the resulting codebase and the author's own direction.

After using this tool, the author reviewed and edited the content as needed and takes full
responsibility for the content of this report and of the underlying project.

\tableofcontents

\section{Concept}\label{concept}

Orbital is a set of web services (a gateway and two backend services) fronting a shared
persistence/caching layer, paired with a Kotlin Multiplatform client application (desktop and
web) that consumes them.

\paragraph{Use case.} The use case is a personal market-data and news dashboard: a user wants to
see current cryptocurrency prices, a coin's price history, and recent news, without visiting
several separate sites.

\begin{itemize}
  \item \textbf{What:} The software continuously ingests cryptocurrency prices and news, caches
  and rate-limits access to that data at a single gateway, and presents it through a searchable,
  filterable client.
  \item \textbf{Roles:} There is currently a single, anonymous role, since there are no user
  accounts and nothing distinguishes one caller from another beyond their IP address for
  rate-limiting purposes (see \cref{security} for a proposed design that would introduce real
  roles).
  \item \textbf{Interaction:} Users interact via a desktop window or a web browser. Interaction is
  pull-based and relatively low-frequency from any single user's perspective (opening the app,
  searching for a coin, expanding its chart, switching a news category), but the client also
  auto-refreshes periodically, and many independent users may be doing this concurrently.
  \item \textbf{Data Storage:} The system does not store any user-specific data, since there is no
  concept of a user account to attach data to. The data it \emph{does} store, price snapshots and
  news articles, is global, shared across every user, persisted in two separate PostgreSQL
  databases (one per backend service).
\end{itemize}

\paragraph{Why is distribution needed?} Distribution was deliberately chosen for the following
reasons, not imposed artificially:

\begin{itemize}
  \item \textbf{Resource Sharing:} All clients observe the same market and news data, so it makes
  sense to fetch and store it once, centrally, rather than have every client poll CoinGecko
  independently, which would also multiply load against CoinGecko's own rate
  limits~\citep{coingecko2024api}.
  \item \textbf{Fault Tolerance:} A single unreplicated database or cache instance would make the
  whole platform's persistence path a single point of failure; \cref{fault-tolerance} describes
  how this is addressed.
  \item \textbf{Independent Scalability:} The gateway is the one component every client talks to,
  so it is the one component that needs to be horizontally scalable independently of the two
  backend services behind it; \cref{availability} describes how this was designed for.
  \item \textbf{Educational Purpose:} To implement and analyze real distributed-systems patterns,
  namely streaming database replication, quorum-based cache failover, a distributed rate limiter,
  and CAP-theorem trade-offs, in a project small enough to reason about end to end.
\end{itemize}

\section{Requirements Elicitation and Analysis}\label{requirements}

\subsection{Functional Requirements}\label{functional-requirements}

\begin{description}
  \item[FR1] The system shall provide current market data (price, market cap, 24h volume, 24h
  price change) for the top $N$ cryptocurrencies by market capitalization. \\
  \emph{Acceptance:} \texttt{GET /api/v1/market/prices} returns a JSON array of at least $N$
  coins, each carrying all of the above fields.

  \item[FR2] The system shall provide historical price data for an individual cryptocurrency,
  built from the system's own continuously-ingested observations rather than only the upstream
  API's short-range sample. \\
  \emph{Acceptance:} \texttt{GET /api/v1/market/history/\{coinId\}} returns a chronologically
  ordered series of prices for a coin that has previously been observed by the system.

  \item[FR3] The system shall provide a list of news articles, optionally filtered by category. \\
  \emph{Acceptance:} \texttt{GET /api/v1/news[?category=]} returns articles consistent with the
  requested category filter, or all categories if none is given.

  \item[FR4] The system shall provide read access to its own previously-stored news archive,
  independent of the live upstream feed's availability. \\
  \emph{Acceptance:} \texttt{GET /api/v1/news/history[?category=]} returns previously persisted
  articles even when the upstream news source is unreachable.

  \item[FR5] The client application shall allow a user to search and filter cryptocurrencies by
  free text, and news articles by category. \\
  \emph{Acceptance:} typing into the market search field narrows the displayed coin list;
  selecting a news category narrows the displayed articles to that category.

  \item[FR6] The client application shall present an individual coin's historical price trend
  visually. \\
  \emph{Acceptance:} expanding a coin in the client renders a chart built from FR2's data,
  falling back to the live snapshot's short-range sample if no history has been ingested yet.

  \item[FR7] The system shall continuously and automatically ingest market data on a fixed
  interval, independent of API request traffic, and shall perform a one-time historical backfill
  for each newly observed coin. \\
  \emph{Acceptance:} price snapshots keep accumulating in storage over time even with zero API
  traffic; a coin observed for the first time has a long history available immediately rather
  than only from the moment it was first observed.
\end{description}

\subsection{Non-functional Requirements}\label{non-functional-requirements}

\begin{description}
  \item[NFR1 (Availability)] The gateway shall keep serving already-cached responses even if a
  backend service, the cache, or the rate limiter's backing store is temporarily unavailable,
  rather than failing the request outright. \\
  \emph{Acceptance:} a Redis outage must not cause the gateway to reject requests outright
  (verified in this project by the rate limiter's fail-open behaviour, \cref{fault-tolerance}).

  \item[NFR2 (Fault Tolerance)] Loss of a single PostgreSQL primary shall not lose previously
  committed data, given a configured streaming-replication standby. \\
  \emph{Acceptance:} \texttt{pg\_stat\_replication} on a primary shows an actively streaming
  replica with a matching log position (verified manually in this project, \cref{acceptance-test}).

  \item[NFR3 (Scalability)] The gateway shall be horizontally scalable to multiple replicas
  without each replica enforcing an independent, and therefore effectively multiplied, rate
  limit. \\
  \emph{Acceptance:} a shared, Redis-backed counter enforces one limit per client across every
  replica (implemented and unit-tested; full multi-replica behaviour under load is discussed as
  not yet verified in \cref{future-works}).

  \item[NFR4 (Performance)] Market and news read endpoints requested repeatedly within a short
  window shall be served from cache rather than re-fetching upstream on every request. \\
  \emph{Acceptance:} two identical requests issued within the cache TTL result in exactly one
  upstream call (covered by an automated test, \cref{automatic-testing}).

  \item[NFR5 (Observability)] All backend services shall expose health and metrics endpoints,
  scraped and visualized centrally. \\
  \emph{Acceptance:} \texttt{/health} and \texttt{/metrics} respond on the gateway and both
  backend services; a Prometheus/Grafana dashboard is provisioned automatically with no manual
  setup.
\end{description}

\subsection{Implementation Choices}\label{implementation-choices}

\begin{enumerate}
  \item \textbf{Technical Stack:} The entire codebase, backend, client, and shared library
  modules, is written in Kotlin, with Ktor~\citep{ktor2024docs} as the backend HTTP
  server/client framework.
  %
  \begin{itemize}
    \item \emph{Justification:} Ktor's coroutine-based concurrency model suits the I/O-bound
    proxy/aggregation workload of a gateway and two backend services well, and keeping the whole
    backend on one language and build toolchain removes a language boundary that would otherwise
    need to be crossed for no architectural reason.
  \end{itemize}

  \item \textbf{Client Platform:} Kotlin Multiplatform with Compose Multiplatform, targeting both
  JVM desktop and \texttt{wasmJs} web from a single codebase.
  %
  \begin{itemize}
    \item \emph{Justification:} avoids maintaining two separate UI codebases for what is
    otherwise the same client logic, and keeps the client in the same language as the backend.
  \end{itemize}

  \item \textbf{Database:} PostgreSQL 15~\citep{postgresql2024docs}, accessed through the Exposed
  ORM, with one independent database per backend service (\texttt{market\_data},
  \texttt{news\_data}).
  %
  \begin{itemize}
    \item \emph{Justification:} a mature, well-understood relational engine with first-class
    streaming replication, directly relevant to NFR2. Separate databases per service keep each
    service's persistence failure domain isolated from the other's.
  \end{itemize}

  \item \textbf{Cache and Coordination:} Redis, accessed via Jedis, behind either a plain
  \texttt{JedisPool} or a Sentinel-aware \texttt{JedisSentinelPool}~\citep{redis2024sentinel}
  depending on configuration.
  %
  \begin{itemize}
    \item \emph{Justification:} Redis is the de facto standard for this kind of short-TTL proxy
    cache, and Sentinel is its own built-in high-availability mechanism rather than a
    third-party add-on, keeping the operational surface small.
  \end{itemize}

  \item \textbf{Load Balancing:} nginx, in front of however many gateway replicas are running.
  %
  \begin{itemize}
    \item \emph{Justification:} a minimal, well-understood reverse proxy was sufficient; nothing
    in this project's requirements called for a full service mesh.
  \end{itemize}

  \item \textbf{Deployment:} Docker Compose~\citep{docker2024compose} orchestrating the entire
  stack on a single host.
  %
  \begin{itemize}
    \item \emph{Justification:} reproducibility across development and this report's own manual
    verification work (\cref{acceptance-test}), without the operational overhead of a
    multi-host orchestrator that this project's scale does not need.
  \end{itemize}
\end{enumerate}

\subsection{Relevant Distributed System Features}\label{ds-features}

\paragraph{Transparency.} Real, and central to the design. \texttt{orbital-app} only ever talks
to the gateway (\cref{architecture}); it has no knowledge of there being two backend services,
two PostgreSQL primaries, a Redis Sentinel quorum, or (potentially) several gateway replicas
behind a load balancer. Location and replication are meant to be invisible to it by construction.

\paragraph{Fault tolerance, dependability.} Highly relevant, and the area with the most concrete
work in this project: PostgreSQL streaming replication for both databases, a Redis Sentinel
quorum, retry-with-backoff on startup, and fail-open behaviour in both the cache and the rate
limiter. Uninterrupted service is not strictly required (this is not a payments system), but
losing previously-collected price/news history to a single container crash was judged
unacceptable given FR2/FR4's whole point is to keep that history around. \Cref{fault-tolerance}
details how.

\paragraph{Scalability.} Relevant specifically at the gateway, the one component every client
and every request passes through; not particularly relevant at the two backend services, which
are not designed to need more than one replica each. \Cref{availability} covers the design;
\cref{future-works} is explicit about what has and has not actually been exercised under load.

\paragraph{Security, trust.} The weakest feature in the current implementation (see
\cref{security}): there are no user accounts, so technically there is no \emph{sensitive} data to
protect yet, and no authentication or authorization exists. This is acceptable for a public
read-only aggregator, but it is a real gap relative to any future version of the product that
introduces accounts, which is why \cref{security} proposes a concrete design rather than leaving
the question unaddressed.

\paragraph{Resource sharing.} Central to the system's existence (\cref{concept}): price and news
data is fetched once and shared across every client, rather than each client fetching
independently; the Redis cache, the Redis-backed rate-limit counters, and both PostgreSQL
databases are all resources shared across however many gateway/service replicas are running.

\paragraph{Openness, interoperability.} Relevant at the boundaries: the backend talks to
CoinGecko's public REST/JSON API as an external, independently-versioned system, and the client
and backend communicate over plain HTTP/JSON rather than a project-specific binary protocol, so
either side could in principle be replaced independently as long as that contract holds.

\paragraph{Evolvability, maintainability.} Actively designed for: the backend is split into
independently buildable Gradle modules, dependency versions are centralized, and every push runs
automated tests, static analysis (detekt), and formatting checks (ktfmt/spotless) in CI, with
releases cut automatically from Conventional Commit history. This matters because the project
has already gone through several rounds of adding infrastructure (replication, Sentinel, a
distributed rate limiter) without a rewrite.

\paragraph{Performance, concurrency.} Relevant but not under strict quantitative targets: there
is no SLA on response time. Ktor's coroutine-based I/O and the proxy cache (\cref{availability})
keep the gateway from doing redundant work, which matters more for not wasting CoinGecko's own
rate-limited quota than for raw latency.

\paragraph{Economy, costs.} A real, concrete driver of at least one implementation decision
already made: news-service currently serves a mocked fallback feed rather than paying for a
NewsAPI subscription, a cost-based trade-off made explicitly rather than left implicit.

\section{Design}\label{design}

This chapter explains the strategies used to meet the requirements identified in
\cref{requirements}, motivated against the features identified in \cref{ds-features}.

\subsection{Architecture}\label{architecture}

Orbital uses a \textbf{Microservices} architecture decomposed along functional boundaries. After
analysis, three macro concerns were identified:

\begin{itemize}
  \item \textbf{Market Data Logic:} Performed by \texttt{market-service} and the standalone
  \texttt{market-ingestor} worker. Handles CoinGecko integration, price persistence, and
  historical backfill. Exposes a REST API for reads; the worker itself exposes nothing and only
  writes.
  \item \textbf{News Logic:} Performed by \texttt{news-service}. Handles the (currently mocked)
  news feed and its own persisted archive. Exposes a REST API.
  \item \textbf{Edge:} Handled by \texttt{gateway} (the single client-facing entry point, caching
  and rate limiting) and \texttt{gateway-lb} (the load balancer in front of however many gateway
  replicas are running).
\end{itemize}

This \textbf{functional decomposition} ensures that market-service and news-service, two
independent external integrations with independent failure modes, can evolve, be deployed, and
fail independently of one another, and that the one component actually needing to scale (the
gateway) is not coupled to the two that do not. The client, \texttt{orbital-app}, is
architecturally a classic thin client: all domain logic and persistence live behind the gateway,
and the client's own role is presentation and local UI state only.

\subsection{Infrastructure}\label{infrastructure}

The infrastructure is composed of the following decoupled components, all orchestrated by one
\texttt{docker-compose.yml}:

\begin{itemize}
  \item \textbf{Client-side:} \texttt{orbital-app}, a Compose Multiplatform application (desktop
  JVM and \texttt{wasmJs} web). It communicates with the gateway exclusively via REST/HTTP.
  \item \textbf{Edge Layer:} \texttt{gateway-lb}, an nginx instance acting as a reverse proxy. It
  provides a single entry point on port \texttt{8080}, hides however many \texttt{gateway}
  replicas sit behind it, and re-resolves the \texttt{gateway} service name on every request
  rather than caching it once, so scaling is actually reflected.
  \item \textbf{Service Layer:} \texttt{gateway} (routing, caching, rate limiting),
  \texttt{market-service} and \texttt{news-service} (independent Ktor applications, each owning
  its own database), and \texttt{market-ingestor} (a worker process with no HTTP interface).
  \item \textbf{Persistence Layer:} Two independent PostgreSQL 15 primaries:
  \texttt{postgres-market} (owned by market-service/market-ingestor) and \texttt{postgres-news}
  (owned by news-service), each with its own streaming-replication standby.
  \item \textbf{Cache / Coordination Layer:} A Redis primary and replica, monitored by a
  three-node Sentinel quorum.
  \item \textbf{Monitoring:} Prometheus (scraping all three HTTP services) and Grafana.
  \item \textbf{Administration:} pgAdmin, for inspecting both PostgreSQL instances.
\end{itemize}

\begin{quote}
\textit{[A component diagram of this topology, covering the client, gateway-lb, gateway, the two
backend services, the ingestor, both database pairs, the Redis/Sentinel cluster, and monitoring,
would be inserted here.]}
\end{quote}

All of these currently run on a single host, on one Docker bridge network
(\texttt{orbital-net}); nothing in the design assumes same-host placement, but nothing in the
current deployment spans more than one host either. That would be the natural next step if
Orbital needed genuine geographic distribution, which \cref{concept} does not identify as a
driver for this product.

Components find each other by \textbf{name}, not by hardcoded IP address: Docker Compose's
embedded DNS resolves each service's name (e.g.\ \texttt{postgres-market}, \texttt{redis},
\texttt{redis-sentinel-1}) to whichever container is currently running it. For Redis specifically,
this is deliberately \emph{not} enough on its own: a client that resolved \texttt{redis} once and
cached the connection would keep talking to the old master after a Sentinel failover. Both the
gateway and news-service therefore talk to Redis through the Sentinel quorum
(\texttt{REDIS\_SENTINEL\_HOSTS}) rather than a fixed hostname, asking the quorum who the current
master is on every new connection. This makes the infrastructure \textbf{location-transparent} in
the sense relevant to \cref{ds-features}: nothing downstream of the gateway needs to know which
physical Redis node is currently serving as master.

\subsection{Modelling}\label{modelling}

\subsubsection{Domain Entities}\label{modelling-entities}

The domain is partitioned between the two backend services, each owning its database outright:

\begin{itemize}
  \item \textbf{Price Snapshot (market-service):} coin id, symbol, name, current price, market
  cap, 24h volume, 24h change, and an observation timestamp. One row per observation.
  \item \textbf{News Article (news-service):} id, title, description, URL, source, category, and
  a published timestamp. One row per observation.
\end{itemize}

Both tables map to exactly one infrastructural location each: price snapshots live only in
\texttt{postgres-market}'s \texttt{market\_prices} table, news articles only in
\texttt{postgres-news}'s \texttt{news} table. There is no replication of domain data
\emph{between} the two backend services, only within each one's own PostgreSQL primary/standby
pair.

\begin{quote}
\textit{[An entity diagram of \texttt{market\_prices} and \texttt{news} would be inserted here.]}
\end{quote}

\subsubsection{The Backfill Gate}\label{modelling-backfill-gate}

A single idempotency check, \texttt{MarketRepository.hasHistory(coinId)}, is the mechanism that
makes FR7's one-time historical backfill safe to re-run unconditionally on every startup of
\texttt{market-ingestor}, rather than needing its own separate completion-tracking table. Before
fetching a coin's full \texttt{market\_chart} history, the worker checks whether any row already
exists for that coin; if one does, the backfill is skipped entirely, with no CoinGecko call and
no delay. This single boolean gate is what lets the worker's startup sequence be simply ``for
every top coin, backfill if needed, then enter the polling loop'' without the worker needing to
remember, across restarts, which coins it has already handled: that memory already exists, as a
side effect, in the very table the backfill writes to. This same idempotency is what makes the
worker's restart behaviour (\cref{fault-tolerance}) safe, since a crash mid-backfill simply means
the next startup re-checks the gate and continues from wherever the data actually stopped.

There is no equivalent gate on the regular polling loop, by design: each polling tick is expected
to append a fresh row regardless of what came before, since FR7's whole point is a continuously
growing observation log, not a single current-value row to be kept idempotent against repeats.

\subsection{Interaction}\label{interaction}

All interaction in Orbital is synchronous \textbf{request/response} over HTTP. The client
(\texttt{orbital-app}) issues a request to the gateway; the gateway either answers directly from
its Redis/in-process cache or proxies the request onward to the relevant backend service, which
in turn either answers from its own database or makes an outbound HTTP call to CoinGecko (or the
mocked news source) before persisting the result and responding. There is exactly one interaction
pattern in use throughout, client-server request/reply, chosen deliberately over, for instance, a
publish/subscribe or streaming pattern, since nothing in \cref{requirements} calls for push
updates: the client's auto-refresh (FR5/FR6) is itself just a periodic pull, not a
server-initiated push.

\begin{quote}
\textit{[A sequence diagram of a cached vs.\ uncached \texttt{GET /api/v1/market/prices} request,
client $\to$ gateway $\to$ (cache hit, or market-service $\to$ CoinGecko), would be inserted
here.]}
\end{quote}

\subsection{Behaviour}\label{behaviour}

The gateway and both backend services are \textbf{stateless} between requests: any state they
touch (the cache, the rate-limit counters, the databases) lives outside the process, in Redis or
PostgreSQL, which is precisely what makes them safe to run as multiple replicas
(\cref{availability}). The market-ingestion worker is the one genuinely \textbf{stateful}
component: it holds an in-process polling loop with its own timing, and is not designed to be run
as more than one replica; doing so would duplicate every fetch and every write, since nothing
coordinates multiple ingestor instances against each other (this constraint is discussed further,
as a known limitation, in \cref{future-works}).

State updates therefore happen in exactly two places: the ingestion worker appends price
snapshots on its fixed interval (FR7), and each backend service appends whatever it fetched
immediately after answering a live (non-cached) request to it.

\subsection{Data and Consistency Issues}\label{data-and-consistency-issues}

\subsubsection{CAP Theorem Position}\label{cap-theorem-position}

The system deliberately prioritizes \textbf{Availability over Consistency} (AP), the opposite
choice from a transactional system, and the correct one here: a client being served a price that
is a few seconds stale is an acceptable, routine outcome; a client being refused an answer
because the authoritative value could not be confirmed is not, for a read-only aggregator with no
financial transactions riding on the exact figure. Framed against the CAP
theorem~\citep{brewer2000towards,gilbert2002brewer}, under a network partition or a dependency
outage the system chooses to keep answering, whether from cache or by falling open past a failed
rate-limit check, rather than to block or reject.

This AP choice is why both the gateway's cache and its distributed rate limiter are built to
\textbf{fail open} (\cref{fault-tolerance}) rather than fail closed, and why PostgreSQL
replication (\cref{fault-tolerance-replication}) is asynchronous: a CP alternative was considered
and rejected, since making every write wait for synchronous standby acknowledgement would trade
away write availability for a consistency guarantee (zero data loss on failover) that nothing in
\cref{requirements} actually demands.

\subsubsection{Consistency Model}\label{consistency-model}

Both backend services persist every fetch to PostgreSQL, via the Exposed ORM over a
HikariCP-managed connection pool, as \textbf{append-only} time series: a repeated observation of
the same coin or the same article produces a \emph{new} row rather than overwriting an existing
one. This is a deliberate modelling decision (\cref{modelling}), not an oversight: it means
\texttt{getPriceHistory}/\texttt{getHistory} read a genuine time series rather than a single
current value, which is exactly what FR2/FR4 need, at the cost of the tables only ever growing
(mitigated by an index on \texttt{(coin\_id, last\_updated)} / \texttt{(category,
published\_at)} and a capped result size on every read). Reads of the cache, meanwhile, are only
ever eventually consistent with the underlying database within the cache's short TTL, an
explicit, accepted trade-off given the AP position above.

\subsubsection{Shared Data and Synchronization}\label{shared-data-sync}

No domain data is shared \emph{between} market-service and news-service: each owns its own
database outright, and neither ever queries or writes to the other's. Because writes to each
table are append-only, concurrent writers never contend on the same row; there is no update lock,
no optimistic-concurrency version column, and no need for one, since nothing is ever updated in
place. The one piece of state genuinely shared across component boundaries is the Redis cache and
the Redis-backed rate-limit counters, shared across however many gateway (and, for the cache,
news-service) replicas are running, which is the entire reason they live in Redis rather than in
each process's own memory (\cref{availability}).

\subsection{Fault-Tolerance}\label{fault-tolerance}

\subsubsection{Service Isolation via Independent Databases}\label{fault-tolerance-isolation}

Because market-service and news-service are independent processes with independent databases and
independent HikariCP connection pools, the failure or overload of one cannot directly cause the
other to fail. A connection-pool saturation in news-service, for instance, does not affect
market-service's pool.

\subsubsection{Replication}\label{fault-tolerance-replication}

Both PostgreSQL primaries (\texttt{postgres-market}, \texttt{postgres-news}) run with a dedicated
streaming-replication standby~\citep{postgresql2024docs}, set up via a replication role and
\texttt{pg\_basebackup}. This exists purely for durability: losing a primary's container and
volume would otherwise lose every price snapshot and news article ever collected, directly
undermining FR2/FR4/NFR2. As discussed in \cref{cap-theorem-position}, replication here is
asynchronous by deliberate choice; the practical consequence, honestly stated, is that a primary
failing \emph{at the exact moment of a write} could still lose that one write before it reached
the standby, an accepted trade-off given this project's AP position, not an oversight. Redis
similarly runs a primary/replica pair, monitored by a three-node Sentinel quorum (quorum size 2)
that would promote the replica on a primary failure~\citep{redis2024sentinel}.

\subsubsection{Fail-Open Cache and Rate Limiter}\label{fault-tolerance-fail-open}

Both the gateway's proxy cache and its distributed rate limiter are designed to \textbf{fail
open}: if a Redis/Sentinel read throws, the gateway logs a warning and falls through to the
upstream fetch (for the cache) or allows the request through anyway (for the rate limiter),
rather than failing the whole request. This was not the original behaviour of the rate limiter:
an internal audit of the system found that it originally had no such fallback and would have
taken down \emph{all} gateway traffic on a Redis blip; this was fixed before this report was
written. Both backend services similarly catch and log (rather than propagate) a per-row
persistence failure during a batch save, so one malformed row cannot take an entire successful
fetch down with it.

\subsubsection{Idempotent Backfill on Restart}\label{fault-tolerance-backfill}

As described in \cref{modelling-backfill-gate}, the \texttt{hasHistory} check that gates FR7's
backfill is what makes \texttt{market-ingestor}'s crash recovery trivial: on restart, the worker
simply re-runs its startup sequence against every top coin, and the gate itself determines which
coins (if any) still need backfilling. No action log, checkpoint file, or separate ``which coins
are done'' table is needed; the already-persisted data \emph{is} the checkpoint.

\subsubsection{Database Connection Retry and Startup Ordering}\label{fault-tolerance-retry}

The market-ingestion worker retries its initial database connection with a capped linear backoff
rather than crashing on a transient startup race against PostgreSQL still initializing. For this
one component, failing to even start is worse than waiting, since the worker exists solely to
write to the database. At the orchestration level, every Docker Compose service declares a
\texttt{healthcheck}, and dependent services wait on \texttt{condition: service\_healthy} rather
than merely on the container having started, so a service never attempts to use a dependency that
has launched but is not yet actually ready.

\subsection{Availability}\label{availability}

\subsubsection{Caching}\label{availability-caching}

The gateway caches upstream responses (market prices, market history, news, news history) keyed
by the request's distinguishing parameters, with a short TTL. When Redis is configured, this
cache is shared across every gateway replica, which is what actually makes running more than one
replica worthwhile, since otherwise each replica would cache (and warm its cache) independently.
When Redis is not configured, the gateway falls back to a bounded in-process cache, which is
correct for a single instance but would silently stop being a \emph{shared} cache across multiple
replicas, a limitation this project's design is explicit about rather than hiding.

\subsubsection{Load Balancing}\label{availability-load-balancing}

An nginx instance (\texttt{gateway-lb}) fronts however many \texttt{gateway} containers are
running, forwarding to the \texttt{gateway} service name with DNS re-resolution on every request
rather than once at startup, so that scaling the number of gateway replicas up or down is
actually reflected. This is paired with the Redis-backed distributed rate limiter
(\cref{fault-tolerance}): load balancing across replicas would otherwise multiply the effective
rate limit by the number of replicas, since Ktor's own built-in rate limiter keeps its counters
in-process.

\subsubsection{Network Partition Behaviour}\label{availability-partitions}

Two different partition scenarios are handled differently by design. A partition that isolates a
PostgreSQL standby from its primary simply stalls replication until it heals, since the primary
keeps serving reads and writes unaffected, given that replication here is asynchronous and
one-way. A partition that isolates a minority of the Redis Sentinel quorum (at most one of three
Sentinels, given quorum size 2) does not block failover decisions, since the remaining two
Sentinels still form a quorum; a partition that isolated a \emph{majority} of Sentinels would
correctly block failover rather than risk a split-brain promotion, at the cost of not being able
to fail over at all until it heals. That is a standard, intentional trade-off of quorum-based
systems~\citep{redis2024sentinel}, not a flaw specific to this deployment.

\subsection{Security}\label{security}

\subsubsection{Current State}\label{security-current}

Orbital currently implements neither authentication nor authorization: every endpoint exposed by
the gateway (\texttt{/api/v1/market/*}, \texttt{/api/v1/news/*}) is publicly readable, and no
cryptographic token scheme protects any of them. This is a deliberate simplification given the
product's current scope, a public, read-only market-data and news aggregator with no user
accounts and no private data, but it does leave the security dimension of \cref{ds-features}
largely unaddressed by the implementation as it stands. The remaining subsections of this section
describe a concrete design for closing that gap; \textbf{none of it has been implemented}: no
Keycloak container exists in the project's \texttt{docker-compose.yml}, the gateway installs no
authentication plugin, and \texttt{orbital-app} has no login flow. It is presented as a proposal,
and listed again as a concrete next step in \cref{future-works}.

\subsubsection{Proposed Identity Design (Not Implemented)}\label{security-identity}

Rather than have each service manage credentials independently, Orbital would introduce a
dedicated identity provider, Keycloak, as a new infrastructural component alongside the existing
Postgres/Redis/Prometheus stack (\cref{infrastructure}). Keycloak would own a single realm for
the project and issue signed JSON Web Tokens (JWTs,~\citep{jones2015jwt}) after an OpenID Connect
(OIDC)~\citep{sakimura2014oidc} Authorization Code flow with PKCE, the flow recommended for
public clients that cannot keep a client secret, which matches \texttt{orbital-app}'s nature as a
Kotlin Multiplatform client distributed as a desktop binary and a static web build, neither of
which can safely embed one.

Centralizing identity this way is a natural fit for Orbital's existing architecture rather than a
bolt-on. The client already talks exclusively to the gateway (\cref{architecture}), so the
gateway would remain the single trust boundary where tokens are validated, exactly as it is
already the single point where caching and rate limiting are enforced. Validation itself would be
stateless: the gateway would install an \texttt{Authentication} plugin with a JWT provider,
fetching Keycloak's signing keys from its JWKS endpoint and caching them locally, refreshed on a
TTL, the same pattern already used elsewhere in the system for distributed state
(\cref{availability}). Because JWT validation needs no round-trip to Keycloak nor a shared session
store, it composes cleanly with the horizontal-scalability work already done on the gateway: any
number of gateway replicas behind the load balancer could validate tokens independently, with no
new shared state and no new single point of failure introduced by the validation step itself.

\subsubsection{Proposed Authorization Model (Not Implemented)}\label{security-authz}

Authorization would be expressed through realm roles embedded as claims in the JWT (e.g.\
\texttt{user}, \texttt{premium}, \texttt{admin}), rather than through a separate access-control
service. Given Orbital's current product shape, the most natural use of these roles would not be
gatekeeping the existing read endpoints (which are deliberately public) but tiering them:
authenticated requests could be granted a materially higher entry in the distributed rate limiter
(\cref{availability}) than anonymous ones, and any future write-capable or administrative endpoint
(e.g.\ forcing a cache invalidation, or re-triggering a historical backfill in the market
ingestion worker) would require the \texttt{admin} role. This keeps the public-aggregator
character of the product intact while still giving the system a real authorization story.

\subsubsection{Keycloak's Own Fault Tolerance (Not Implemented)}\label{security-keycloak-ha}

Keycloak itself would need the same fault-tolerance treatment already applied to Orbital's other
stateful components (\cref{fault-tolerance}): it requires its own relational database for realm
and user state, and running it as a single instance would reintroduce exactly the kind of single
point of failure that the Postgres and Redis replication work in this project was meant to
eliminate. A production-oriented deployment would therefore either run Keycloak against one of the
already-replicated Postgres instances, on a dedicated schema, or, more robustly, run Keycloak
itself in clustered mode across two or more nodes.

\section{Implementation}\label{implementation}

\subsection{Network Protocols and Data Representation}\label{impl-network}

Every component communicates over \textbf{HTTP/1.1}, carrying \textbf{JSON} bodies (via
\texttt{kotlinx.serialization}): client to gateway, gateway to backend services, and backend
services to CoinGecko. No binary or RPC protocol was introduced anywhere in the system; HTTP/JSON
was sufficient for the synchronous request/reply pattern identified in \cref{interaction}, and
keeps every hop debuggable with a plain \texttt{curl}, which mattered in practice during this
project's own manual verification work (\cref{acceptance-test}).

\subsection{Data Persistence and Querying}\label{impl-persistence}

Databases are queried with plain \textbf{SQL}, through Exposed's typed query DSL rather than raw
query strings, which also means no part of this project builds a SQL string by hand from request
input, closing off SQL injection as an attack surface by construction rather than by discipline.
Schema is applied automatically at service startup via Exposed's
\texttt{createMissingTablesAndColumns}, which creates missing tables, columns, and indexes,
verified empirically in this project to also retroactively add a newly-introduced index to an
already-existing table (\cref{automatic-testing}): the scenario an already-running deployment
being upgraded would actually hit.

\subsection{Authentication and Authorization}\label{impl-auth}

As detailed in \cref{security}, no component is currently authenticated or authorized; there is
no OAuth/JWT validation and no RBAC/ABAC check anywhere in the running system. A concrete,
unimplemented design for closing this gap (OIDC via Keycloak, JWTs validated at the gateway, role
claims driving a tiered rate limit) is proposed in \cref{security} and listed again in
\cref{future-works}.

\subsection{Frontend Implementation}\label{impl-frontend}

\texttt{orbital-app} is built with Compose Multiplatform, sharing one UI codebase across a
\texttt{commonMain} source set and platform-specific \texttt{desktopMain}/\texttt{wasmJsMain}
targets. State management uses Compose's own primitives directly (\texttt{remember},
\texttt{mutableStateOf}, and \texttt{LaunchedEffect}) rather than an external state-management
library; for instance, a coin row's expanded/collapsed state and its lazily-fetched price history
are both plain \texttt{remember}ed state, with \texttt{LaunchedEffect(coin.coinId)} triggering
the history fetch only once a coin is actually expanded, and falling back to the live snapshot's
short-range \texttt{sparklineIn7d} data if the fetch returns fewer than two points (FR6). An
injectable \texttt{fetchPriceHistory} function parameter is threaded down through the composable
tree specifically so the chart-fetching behaviour can be swapped out or omitted (as the
backend-free demo build does) without changing the UI components themselves. Network access goes
through a thin \texttt{MarketApiClient} wrapping a Ktor HTTP client, calling the gateway's
\texttt{/api/v1/market/prices} and \texttt{/api/v1/market/history/\{coinId\}} endpoints directly;
the client never talks to market-service or news-service itself (\cref{architecture}). Free-text
coin search and news category filtering (FR5) are both plain client-side list filtering over
whatever the gateway most recently returned, not separate backend endpoints.

\subsection{Containerization and Orchestration}\label{impl-containers}

Every backend component is containerized with Docker. Each Ktor service (gateway, market-service,
news-service) has a multi-stage \texttt{Dockerfile}: a build stage using
\texttt{eclipse-temurin:21-jdk} to produce a fat JAR via Gradle's \texttt{shadowJar}, and a
runtime stage using the much smaller \texttt{eclipse-temurin:21-jre}. \texttt{market-ingestor}
reuses market-service's own build output via a separate \texttt{Dockerfile.ingestor} that differs
only in its entrypoint (\texttt{com.orbital.market.MarketIngestorKt} instead of the HTTP server's
\texttt{main}), rather than duplicating the build. One \texttt{docker-compose.yml} defines the
entire system: all 17 services described in \cref{infrastructure}, on one bridge network, with
environment variables for service URLs, credentials, and feature toggles (such as
\texttt{REDIS\_SENTINEL\_HOSTS}) injected per service, and health checks enforcing startup
ordering (\cref{fault-tolerance-retry}). \Cref{deployment} gives the full per-service breakdown.

\subsection{Technological Details}\label{technological-details}

\begin{itemize}
  \item \textbf{Kotlin} across the entire codebase.
  \item \textbf{Ktor}~\citep{ktor2024docs}: every backend HTTP server and client.
  \item \textbf{Compose Multiplatform}: the \texttt{orbital-app} client, JVM desktop and
  \texttt{wasmJs} web from one codebase.
  \item \textbf{Exposed} / \textbf{HikariCP} against \textbf{PostgreSQL
  15}~\citep{postgresql2024docs}: all persistence; \textbf{H2} (PostgreSQL-compatibility mode)
  stands in for it in the automated test suite.
  \item \textbf{Jedis}, behind an internal abstraction supporting either a plain
  \texttt{JedisPool} or a Sentinel-aware \texttt{JedisSentinelPool}~\citep{redis2024sentinel}.
  \item \textbf{nginx}: \texttt{gateway-lb}'s load balancer.
  \item \textbf{Prometheus} and \textbf{Grafana}: metrics collection and visualization, via a
  Micrometer-backed \texttt{/metrics} endpoint on each HTTP service.
  \item \textbf{Docker Compose}~\citep{docker2024compose}: the sole deployment/orchestration
  mechanism.
  \item \textbf{Gradle} (multi-module), with \textbf{detekt} and \textbf{ktfmt}/\textbf{spotless}
  enforced in CI, and \textbf{GitHub Actions} + \textbf{release-please} automating testing,
  Docker image builds, and versioned releases from Conventional Commit history.
\end{itemize}

\section{Validation}\label{validation}

\subsection{Automatic Testing}\label{automatic-testing}

\subsubsection{Test Infrastructure}\label{testing-infrastructure}

Every JVM backend module has its own automated test suite, run with \texttt{./gradlew test} per
module or \texttt{./gradlew check} for the full suite plus static analysis, executed on every
push and pull request against \texttt{main}/\texttt{develop} by the \texttt{Test \& Lint} GitHub
Actions workflow. Persistence code runs against H2 in PostgreSQL-compatibility mode rather than a
mock of the database, exercising the real Exposed queries and real schema. HTTP-level integration
uses Ktor's \texttt{testApplication}, which runs the real routing and plugin pipeline (including
CORS and rate limiting) against an in-process test server, with upstream HTTP dependencies
replaced by Ktor's \texttt{MockEngine}. Pure logic with no external dependency, such as the
distributed rate limiter's allow/reject/fail-open decision, is tested against a hand-written fake
rather than a mocking framework, so the test has no dependency on a real Redis connection ever
being reachable:

\begin{verbatim}
private class FakeRedisOps : RedisOps {
  private val counts = mutableMapOf<String, Long>()
  override fun incrWithExpire(key: String, windowSeconds: Int): Long =
      counts.merge(key, 1L, Long::plus) ?: 1L
}
\end{verbatim}

\subsubsection{Test Coverage}\label{testing-coverage}

Tests are organized per module (20 for market-service, 14 for gateway, 10 for news-service, plus
a handful across the two shared library modules) and cover the following requirement-grouped
scenarios:

\paragraph{Market Data Ingestion and History (FR2, FR7).} Asserts that
\texttt{getPriceHistory} returns a coin's snapshots oldest-first, excluding other coins'
snapshots; that \texttt{hasHistory} is false until a snapshot is stored for that coin and true
afterward (the backfill gate from \cref{modelling-backfill-gate}); and, as a direct empirical
answer to a question raised during this project's own audit, that
\texttt{createMissingTablesAndColumns} actually adds the \texttt{(coin\_id, last\_updated)} index
to a table created \emph{without} it, simulating an already-running deployment being upgraded.

\paragraph{News Persistence and History (FR3, FR4).} Asserts that \texttt{getHistory} returns
the most recently published articles first, optionally filtered by category, and that the
\texttt{/api/v1/news/history} route returns stored articles filtered the same way at the HTTP
level.

\paragraph{Gateway Caching and Rate Limiting (NFR1, NFR3, NFR4).} Asserts that a repeated
request for the same market prices within the cache TTL results in exactly one upstream call
(NFR4); that a cache read failure falls through to the upstream fetch rather than failing the
request (NFR1); that the distributed rate limiter allows calls up to its configured limit and
rejects beyond it, tracking separate client keys independently (NFR3); and that it fails open
when its backing store throws (NFR1, \cref{fault-tolerance-fail-open}), a regression test for a
real bug found and fixed during this project's own internal audit. A further test asserts that a
\texttt{429} response from the rate limiter still carries CORS headers, a regression test for an
ordering bug found during the same audit.

\paragraph{Persistence Error Handling (NFR1-adjacent corner case).} Asserts that a persistence
write violating a column constraint is logged and skipped rather than crashing the surrounding
batch, on both backend services.

\subsubsection{How to Run}\label{testing-how-to-run}

\begin{verbatim}
./gradlew test    # a single module's tests, e.g. ./gradlew :gateway:test
./gradlew check   # full suite, plus detekt and ktfmt/spotless
\end{verbatim}

Test reports are available under each module's \texttt{build/reports/tests/} directory.

\subsection{Acceptance Testing}\label{acceptance-test}

Several properties of this system are either impractical or outright impossible to assert
deterministically in a fast, isolated automated test, because they depend on real multi-process
timing and real network behaviour: replication lag, Sentinel quorum formation and failover
detection, and DNS-based load-balancer discovery. These were instead verified \textbf{manually}
against the real \texttt{docker-compose} deployment:

\begin{itemize}
  \item \textbf{PostgreSQL streaming replication} (NFR2): querying
  \texttt{pg\_stat\_replication} on the \texttt{postgres-market} primary showed an actively
  streaming replica (\texttt{state: streaming}) with matching sent/write/flush/replay log
  positions.
  \item \textbf{Redis Sentinel} (supporting NFR3's shared-cache precondition): a fatal
  configuration error was found and fixed first (Sentinel refuses to resolve a Docker service name
  as a monitor target unless \texttt{resolve-hostnames yes} is set), after which
  \texttt{sentinel master mymaster} against a live Sentinel instance reported the expected
  master, one discovered replica, and the other two Sentinels forming a three-node quorum.
\end{itemize}

\paragraph{Rationale for manual testing.} Both scenarios require a live, multi-container
environment with real inter-process timing (replication catching up, Sentinel's
down-after-milliseconds detection window) that is impractical to assert deterministically in a
fast unit or integration test without a full infrastructure-level testing harness. What was
\emph{not} verified manually, and is reported honestly as outstanding rather than assumed
working, is discussed in \cref{future-works}: the nginx load balancer's actual round-robin
behaviour across multiple scaled \texttt{gateway} replicas, and the distributed rate limiter's
behaviour under real concurrent multi-replica load.

\section{Deployment}\label{deployment}

\subsection{Prerequisites}\label{deployment-prerequisites}

The following software must be installed on the host machine:
\begin{itemize}
  \item \textbf{Docker} with Docker Compose v2.
  \item \textbf{JDK 21} (for running the client/backend directly outside Docker; not required if
  only using \texttt{docker compose}).
\end{itemize}

\subsection{Installation from Scratch}\label{deployment-installation}

\begin{enumerate}
  \item \textbf{Clone the repository and configure environment variables:}
\begin{verbatim}
git clone <repository-url>
cd orbital
cp .env.example .env
\end{verbatim}
  Every variable in \texttt{.env} has a safe, documented default for local development; only
  credential-related variables (\texttt{POSTGRES\_MARKET\_PASSWORD},
  \texttt{POSTGRES\_NEWS\_PASSWORD}, \texttt{POSTGRES\_REPLICATION\_PASSWORD},
  \texttt{GRAFANA\_ADMIN\_PASSWORD}, \texttt{PGADMIN\_ADMIN\_PASSWORD}) must be overridden before
  any deployment reachable over a real network. These replaced plaintext defaults that
  \texttt{docker-compose.yml} previously hardcoded directly, an issue identified and fixed during
  this project's own security audit.

  \item \textbf{Build and start all services:}
\begin{verbatim}
make up        # equivalent to: docker compose up --build -d
\end{verbatim}
  On first startup, each Ktor service applies its schema automatically via Exposed
  (\cref{impl-persistence}) before accepting traffic.

  \item \textbf{Verify deployment:}
\begin{verbatim}
curl localhost:8080/health
curl localhost:8080/api/v1/market/prices
curl localhost:8080/api/v1/news
\end{verbatim}
  Expected outcome: a healthy status from the first call, and live CoinGecko-sourced data from
  the second, within a few seconds.
\end{enumerate}

\subsubsection{Docker Compose Services}\label{deployment-services}

\texttt{docker-compose.yml} defines the following 17 services:

\begin{itemize}
  \item \textbf{gateway:} The API gateway. Proxies, caches, and rate-limits requests to the two
  backend services. No fixed published port: \texttt{gateway-lb} is the host's actual entry
  point, so that it can be scaled to multiple replicas.
  \item \textbf{gateway-lb:} nginx, published on host port \texttt{8080}. Load-balances across
  however many \texttt{gateway} replicas are running.
  \item \textbf{market-service:} The market-data REST API, published on port \texttt{8081}.
  Depends on \texttt{postgres-market} and \texttt{redis}.
  \item \textbf{market-ingestor:} The standalone ingestion worker (FR7). No published port.
  Depends on \texttt{postgres-market}.
  \item \textbf{news-service:} The news REST API, published on port \texttt{8082}. Depends on
  \texttt{postgres-news} and the Redis/Sentinel cluster.
  \item \textbf{postgres-market / postgres-news:} The two independent PostgreSQL 15 primaries,
  published on ports \texttt{5432}/\texttt{5433}. Health-checked via \texttt{pg\_isready} before
  dependent services start.
  \item \textbf{postgres-market-replica / postgres-news-replica:} Streaming-replication standbys
  of the two primaries, cloned via \texttt{pg\_basebackup} on first start.
  \item \textbf{redis:} The Redis cache primary, published on port \texttt{6379}.
  \item \textbf{redis-replica:} A replica of \texttt{redis}, for Sentinel to promote on failover.
  \item \textbf{redis-sentinel-1/2/3:} The three-node Sentinel quorum (quorum size 2) monitoring
  \texttt{redis}/\texttt{redis-replica}.
  \item \textbf{prometheus / grafana:} Metrics collection (port \texttt{9090}) and visualization
  (port \texttt{3000}), with an auto-provisioned ``Orbital Overview'' dashboard.
  \item \textbf{pgadmin:} A PostgreSQL administration UI (port \texttt{5050}) for both databases.
\end{itemize}

All services share a single Docker bridge network (\texttt{orbital-net}), enabling inter-container
communication by service name (\cref{infrastructure}).

\subsubsection{Startup Ordering and Health Checks}\label{deployment-health}

A critical aspect of the deployment is ensuring services start in the correct order. Each service
that has one declares a \texttt{healthcheck}, and dependent services wait on \texttt{condition:
service\_healthy} rather than the weaker \texttt{condition: service\_started}
(\cref{fault-tolerance-retry}): the former waits for the dependency to genuinely be ready to
accept connections, not merely for its container process to have launched. Without this, a
service could start, attempt to connect to a database that has not finished initializing, and
fail. \texttt{market-ingestor} specifically guards against this further with its own
retry-with-backoff (\cref{fault-tolerance-retry}), since it is not gated by Compose health checks
for its database dependency being fully migrated, only for PostgreSQL itself accepting
connections.

\subsection{Rebuilding Individual Services}\label{deployment-rebuild}

A single service can be rebuilt and restarted without affecting the others:
\begin{verbatim}
docker compose up --build gateway
\end{verbatim}

\subsection{Schema Management}\label{deployment-schema}

Unlike a Liquibase/Flyway-based migration tool, Orbital applies schema changes automatically at
each service's startup, via Exposed's \texttt{createMissingTablesAndColumns}
(\cref{impl-persistence}); there is no separate ``run the migrations'' command, and none is
needed: starting (or restarting) \texttt{market-service} or \texttt{news-service} is itself the
migration step. This was a deliberate simplicity trade-off for a project of this size, with the
trade-off made explicit rather than assumed: \cref{automatic-testing} reports an automated test
that specifically verifies this mechanism retroactively adds a newly-introduced index to an
already-existing table, since that guarantee cannot simply be assumed of every Exposed version.

\section{User Guide}\label{user-guide}

Two usage modes exist. The first is direct HTTP access, for scripting or inspection:

\begin{verbatim}
curl localhost:8080/api/v1/market/prices              # live market data
curl localhost:8080/api/v1/market/history/bitcoin      # a coin's ingested price history
curl localhost:8080/api/v1/news?category=crypto        # news, filtered by category
\end{verbatim}

The second, and primary intended use, is \texttt{orbital-app}: a search box and list of
cryptocurrencies with live price/change data; tapping a coin expands it to reveal a price-history
chart (FR6), built from the backend's own ingested history once enough has accumulated, falling
back to the live snapshot's short-range data otherwise; a separate tab lists news articles,
filterable by category; both views auto-refresh periodically without user action. It is started
with:

\begin{verbatim}
./gradlew :orbital-app:run                          # desktop window
./gradlew :orbital-app:wasmJsBrowserDevelopmentRun   # web, dev server + browser tab
\end{verbatim}

A separate, backend-free build of the same client (market data only, fetched directly from
CoinGecko, no gateway required) is deployed automatically on every release to
\url{https://loch3f.github.io/orbital/}.

\begin{quote}
\textit{[Screenshots of the desktop and web client, showing the market list, an expanded coin's
chart, and the news view, would be inserted here.]}
\end{quote}

\section{Self-evaluation}\label{self-evaluation}

\subsection{Aldo Visconti}\label{self-evaluation-aldo}

This is a single-author project, so one section follows rather than one per group member.

\paragraph{Role.} As the sole developer of this project, I was responsible for every aspect of
the system: requirements elicitation, architectural design, backend implementation across the
gateway and both backend services, the market-ingestion worker, the client application, the
infrastructure work described in \cref{fault-tolerance} and \cref{availability}, and both
automated and manual testing.

\paragraph{Strengths.} The decision to decouple market-data ingestion from API traffic into a
standalone worker (FR7, \cref{architecture}) directly solved the problem that motivated it: a
short, API-traffic-dependent price history. Its idempotent backfill gate
(\cref{modelling-backfill-gate}) also turned out to double as the mechanism that makes the
worker's crash recovery essentially free, without a separate checkpointing scheme. The fail-open
design applied to the cache and, after a self-audit caught its absence, to the distributed rate
limiter as well (\cref{fault-tolerance-fail-open}) proved to be the right default for a system
whose CAP position is availability over consistency (\cref{cap-theorem-position}): it is a
materially different, and better, failure mode than the one the rate limiter originally had.
Verifying claims empirically rather than assuming them held up in practice: the Sentinel
configuration had a real, fatal bug (\texttt{resolve-hostnames}) that only surfaced once I
actually brought the quorum up, and the question of whether a schema migration tool would
retroactively index an already-existing table turned out to be worth a real test rather than an
assumption (\cref{automatic-testing}); in both cases, checking rather than assuming caught
something that would otherwise have been a quiet, production-only surprise. Finally, keeping the
backend split into independently buildable, CI-enforced Gradle modules from the start made it
possible to add substantial new infrastructure (replication, Sentinel, a distributed rate
limiter) across several iterations without destabilizing what already worked.

\paragraph{Weaknesses and Limitations.} Working alone meant no second reviewer catching design
blind spots as they happened; the internal audit process I ran against my own work, which is what
actually found the rate limiter's fail-closed bug and the hardcoded database credentials, was a
partial substitute for that, but it is not equivalent to an independent reviewer, and I can't rule
out that it missed things a second person would have caught. More concretely, several gaps are
already honestly documented rather than glossed over: \texttt{orbital-app}, the one component
real users actually touch, has no automated tests at all, unlike every backend module; the
distributed rate limiter and load balancer were built specifically to make horizontal scaling
safe but were never exercised against real concurrent multi-replica load, only against their own
unit tests and a single-instance manual check; there is no authentication or authorization layer,
only a proposed, unimplemented design (\cref{security}); and PostgreSQL failover is manual in
practice, since nothing promotes a standby automatically. \Cref{future-works} lists these and a
few other concrete next steps in more detail.

\section{Future works}\label{future-works}

\paragraph{Automatic database failover.} Both PostgreSQL standbys (\cref{fault-tolerance}) give
Orbital a readable, up-to-date copy of each database, but nothing promotes a standby to primary
or repoints a backend service's \texttt{DB\_URL} automatically if a primary fails. That is
normally the job of a dedicated tool such as Patroni or repmgr, deliberately left out of this
project as a separate, larger piece of work rather than attempted partially.

\paragraph{Verifying scalability under real load.} The nginx load balancer and the distributed
rate limiter were both designed specifically to make running multiple \texttt{gateway} replicas
safe (NFR3), and both compile and pass their own unit/integration tests, but neither has actually
been exercised against several \emph{scaled} replicas under concurrent load, e.g.\
\texttt{docker compose up --scale gateway=3} followed by sustained traffic, checking that the
rate limit really is enforced once, globally, rather than once per replica. This is the natural
next step beyond the manual verification already performed (\cref{acceptance-test}).

\paragraph{Synchronous replication.} \Cref{cap-theorem-position} discusses the AP position taken
for both PostgreSQL pairs; a future iteration could make replication mode configurable per
deployment, trading write latency for a zero-data-loss guarantee on failover where that trade is
justified.

\paragraph{Automated testing of \texttt{orbital-app}.} Every backend module has an automated test
suite (\cref{automatic-testing}); the client application currently has none. This is a real gap
specifically because the client is the one component actual users interact with directly.

\paragraph{Real news integration.} News-service currently serves a mocked fallback feed rather
than a paid NewsAPI subscription (\cref{ds-features}'s economy discussion); wiring up the real
integration behind the same interface it already exposes (\texttt{NewsApiClient}) would be a
contained change once the cost trade-off is revisited.

\paragraph{Authentication and authorization.} As discussed in \cref{security}, Orbital has no
authentication or authorization layer today. The concrete next step is the Keycloak-based design
proposed there: a dedicated identity-provider service issuing OIDC/JWT tokens, validated
statelessly at the gateway, with realm-role claims driving a tiered rate limit between
authenticated and anonymous traffic, and gating any future write-capable endpoint. This remains
future work rather than a completed feature.

\bibliographystyle{plain}
\bibliography{references}

\end{document}
